\section{The collected total at high energy}
\label{sec:comp-total-energy-high}
The complete analytic coefficient gives the high-energy expansion
\begin{equation}
\begin{aligned}
 \sigma_1=\frac{\alpha_{\mathrm T}^3}{m^2}\biggl\{
 \frac1X\biggl[&\frac{\ln^3X}{3}-\frac{\ln^2X}{2}
 +\frac{17}{4}\ln X-\frac{75}{8}-\frac{\pi^2}{2}+4\zeta(3)\biggr]
 +O\!\left(\frac{\ln^3X}{X^2}\right)\biggr\},\quad X\to\infty.
\end{aligned}
\label{eq:comp-collected-total-high}
\end{equation}
Here $X=s_{\mathrm{phys}}/m^2-1$. This expansion is taken after all
virtual and real contributions have been integrated and combined.
The finite constant in the square brackets is determined below from
our GPL primitives, independently of a comparison cross section.

\paragraph{Terms needed from the kernels.} Set $\ell=\ln X$.
Most coefficients in the collected expression are $O(1/X)$, so their
GPLs are needed through the finite constant. Three weight-two
functions multiply coefficients of order one and require an
additional inverse-power term:
\begin{align}
 G(-1,-1;X)&=\frac{\ell^2}{2}+\frac{\ell}{X}
                         +O\!\left(\frac{\ell}{X^2}\right),\\
 G(0,-1;X)&=\frac{\ell^2}{2}+\frac{\pi^2}{6}-\frac1X
                         +O(X^{-2}),\\
 G(r,-1;X)&=\frac{\ell^2}{2}-\frac{\pi^2}{18}
                +\frac{2\ell+1}{X}
                         +O\!\left(\frac{\ell}{X^2}\right).
\end{align}
The first two follow from logarithmic expansion and dilogarithm
inversion. The third follows from \eqref{eq:comp-ir-short-function}
with $q=X^{-1}-2X^{-2}+O(X^{-3})$. Keeping these subleading terms
is essential because the order-one contributions cancel in the sum.

For the six remaining root words, define the constants by
\begin{align}
 G(r,-1,-1;X)&=\ell^3/6+K_a+O(\ell^2/X),\\
 G(r,-1,0;X)&=\ell^3/6-\pi^2\ell/6+K_b+O(\ell^2/X),\\
 G(r,0,-1;X)&=\ell^3/6+\pi^2\ell/6+K_c+O(\ell^2/X),\\
 G(r,-2,-1;X)&=\ell^3/6-\pi^2\ell/12+K_d+O(\ell^2/X),\\
 G(-4,r,-1;X)&=\ell^3/6-\pi^2\ell/18+K_e+O(\ell^2/X),\\
 G(0,r,-1;X)&=\ell^3/6-\pi^2\ell/18+K_f+O(\ell^2/X).
\end{align}
The logarithmic terms follow by integrating their defining kernels at
large $X$. Rational-letter functions follow by the same recursion,
with constants fixed by their classical polylogarithmic expressions.
Before evaluating the six constants, collecting every term gives
\begin{equation}
\begin{aligned}
 C_1(X)=\frac1X\biggl[&\frac{\ell^3}{3}-\frac{\ell^2}{2}
 +\frac{17\ell}{4}-\frac{75}{8}-\frac{\pi^2}{2}+13\zeta(3)\\
 &+15K_a-6K_b-8K_c-16K_d-6K_e-3K_f\biggr]
 +O(\ell^3/X^2).
\end{aligned}
\end{equation}

Using the endpoint constants already evaluated in
\eqref{eq:comp-high-root-constants}, we obtain
\begin{equation}
 15K_a-6K_b-8K_c-16K_d-6K_e-3K_f=-9\zeta(3).
\end{equation}
The Clausen constant cancels exactly. Substitution proves
\eqref{eq:comp-collected-total-high}, including its finite constant.
The rational kernels have inverse-power expansions at infinity, so
the next order is bounded by $O(\ln^3X/X^2)$.

Only now do we compare with Lee, Schwartz and Zhang,
Equation (11)~\cite{qed-processes-lsz}. Replacing $X$ by
$s_{\mathrm{phys}}/m^2$ inside the leading term changes only the displayed
remainder. All three logarithmic coefficients and the finite
coefficient agree with that published expression.
